% =========================================================
% 重大仪器模板
% =========================================================

\documentclass[12pt,UTF8,AutoFakeBold=3,a4paper]{ctexart}

% 1. 加载核心配置
% config.tex
\ProvidesFile{config.tex}[2026 NSFC Common Config]

% =========================================================
% 模块 1：基础宏包引入 
% =========================================================
\usepackage[english]{babel}
\usepackage{xeCJK}
\usepackage{xeCJKfntef}
\usepackage{fontspec} 
\usepackage{pifont}
\usepackage{graphicx}
\usepackage{ulem}
\usepackage{setspace}
\usepackage{enumitem} % 列表控制
\usepackage{amsmath}
\usepackage{mdframed}
\usepackage{xcolor}
\usepackage{amsfonts}
\usepackage{booktabs}
\usepackage[subrefformat=parens,labelformat=parens]{subfig}
\usepackage{hyperref}
\usepackage{titlesec} 
\usepackage{tcolorbox}
\usepackage{amssymb}  % 提供 \blacksquare 等数学符号
\usepackage{bm}       % 提供数学粗体

% =========================================================
% 模块 2：读取外部用户配置 (置于前方，确保全局变量生效)
% =========================================================
% =========================================================
%  NSFC Style Configuration (风格配置文件)
%  说明：此文件仅用于控制排版风格开关，核心逻辑在 config.tex
% =========================================================

% ---------------------------------------------------------
%  开关 1：正文字体设置
% ---------------------------------------------------------
% true  = 正文强制使用【楷体】
% false = 正文使用标准【宋体】
\newif\ifUseKaitiBody
\UseKaitiBodyfalse  % <--- 修改这里：true 或 false

% ---------------------------------------------------------
%  开关 2：标题字体选择
% ---------------------------------------------------------
% 定义各级标题（一级、二级）使用的字体
% 常用选项：\kaishu (楷体), \heiti (黑体), \songti (宋体)
\newcommand{\HeadingFont}{\mysong} 

% ---------------------------------------------------------
%  开关 3：标题编号风格
% ---------------------------------------------------------
% 请输入数字选择一种风格：
% 1 = 标准简洁: 1  / 1.1 / 1.1.1 
% 2 = Word风格: 1. / 1.1 / 1.1.1 (一级标题带个点)
% 3 = 中文风格: (一) / 1. / (1)  (公文风)
\newcommand{\NumberingStyle}{2} % <--- 修改这里：1, 2, 或 3

% ---------------------------------------------------------
%  开关 4：交叉引用与链接颜色
% ---------------------------------------------------------
% 定义参考文献标号、目录、公式引用的颜色。
% 打印版（黑白纸质）强烈建议用 black。电子版可用 blue, red, green, darkgray 等
\newcommand{\ThemeColor}{blue} % <--- 修改这里：输入你想要的颜色单词


% ---------------------------------------------------------
%  开关 5：全局布局与间距调控
% ---------------------------------------------------------
% 设置说明：
% 1. 段间距 (UserParskip): 建议值 0.3ex 到 0.8ex。若篇幅吃紧可设为 0pt。
\newcommand{\UserParskip}{0.5ex} % <--- 修改这里：段间距数值 (如 0.5ex, 3pt)

% 2. 列表模式 (UserListMode): 
%    0 = 极致紧凑: 列表项内外完全不留白 (nosep)，最省空间。
%    1 = 宽松美观: 列表间距随段间距自动按比例缩放，呼吸感强。
\newcommand{\UserListMode}{0}    % <--- 修改这里：0 或 1

%-----------------------------------------------
% 自定义命令（自定义颜色加粗）
%-----------------------------------------------
\definecolor{myemphasize}{RGB}{0, 0, 64} % 经典学术蓝，比纯黑更有质感#000640
% 定义自定义加粗命令 \mybf
\newcommand{\mybf}[1]{{\color{myemphasize}\textbf{#1}}}


% =========================================================
% 模块 3：基础格式与环境配置
% =========================================================
% --- 颜色与路径 ---
\definecolor{NsfcBlue}{RGB}{0,112,192}
\graphicspath{{figures/}}

% --- 参考文献设置 ---
\usepackage[square,numbers,sort&compress,super]{natbib}
\setlength{\bibsep}{1pt plus 0.3ex}
\pagestyle{empty}

% --- 中英文字体与间距设置 ---
\CJKsetecglue{\hskip 0.2em} % 中英文/数字间距
\renewcommand{\baselinestretch}{1.3} 

\setmainfont{Times New Roman}
\setCJKfamilyfont{mysong}[
    UprightFont = {fonts/simsun.ttc}, 
    % 针对粗体单独设置 Scale
    BoldFont    = {fonts/siyuansongti-Bold.otf},
    BoldFeatures = {Scale=0.96} 
]{}

% 创建快捷命令
\newcommand{\mysong}{\CJKfamily{mysong}}

\newfontfamily\KaiEN[Path=fonts/, AutoFakeBold=1.8]{simkai.ttf}
\setCJKfamilyfont{KaiCJK}[Path=fonts/, AutoFakeBold=1.8]{simkai.ttf}
\newcommand{\mykai}{\CJKfamily{KaiCJK}\KaiEN}

% --- 字号快捷命令 ---
\newcommand{\sanhao}{\zihao{3}}
\newcommand{\sihao}{\zihao{4}}
\newcommand{\xiaosihao}{\zihao{-4}}
\newcommand{\wuhao}{\zihao{5}}

% --- 图表标题中文化 ---
\addto\captionsenglish{
    \renewcommand{\refname}{\sihao \mykai \leftline{参考文献}}
    \renewcommand{\tablename}{表}
    \renewcommand{\figurename}{\mykai 图}
}

% --- 列表与代码块环境 ---
\setlist{nosep} % 清除列表多余间距

\newtcolorbox{codebox}{
    colback=gray!5,       
    colframe=gray!30,     
    boxrule=0.5pt,        
    arc=2pt,              
    left=4pt, top=2pt, bottom=2pt, right=4pt, 
    before skip=1ex, after skip=1ex 
}

% --- 超链接颜色设置 ---
\hypersetup{
    colorlinks = true,              
    citecolor  = \ThemeColor,       
    linkcolor  = \ThemeColor,       
    urlcolor   = \ThemeColor        
}

% =========================================================
% 模块 4：标题排版与自定义宏 (NSFC 专用)
% =========================================================

% --- 1. 动态间距长度定义 (底层 Hack 核心) ---
% 预设基础间距
\newlength{\SecTopBase}
\setlength{\SecTopBase}{2.0ex plus .2ex minus .2ex}
\newlength{\SubSecTopBase}
\setlength{\SubSecTopBase}{1.2ex plus .1ex minus .2ex}
\newlength{\SubSubSecTopBase}
\setlength{\SubSubSecTopBase}{0.8ex plus .1ex minus .1ex}

% 定义当前生效的动态间距 (初始等于基础间距)
\newlength{\SecTopCurrent}
\setlength{\SecTopCurrent}{\SecTopBase}
\newlength{\SubSecTopCurrent}
\setlength{\SubSecTopCurrent}{\SubSecTopBase}
\newlength{\SubSubSecTopCurrent}
\setlength{\SubSubSecTopCurrent}{\SubSubSecTopBase}

% --- NSFC 智能标题配置 ---
\newlength{\NsfcHeaderIndent}
\setlength{\NsfcHeaderIndent}{2em}

\newlength{\NsfcSubHeaderIndent}
\setlength{\NsfcSubHeaderIndent}{2em}

\newlength{\NsfcHeaderTopSkip}
\setlength{\NsfcHeaderTopSkip}{0.5em}

\newif\ifIsFirstSubHeader
\IsFirstSubHeadertrue

\newcounter{NsfcHeaderCounter}
\renewcommand{\theHsection}{\arabic{NsfcHeaderCounter}.\arabic{section}}

\newcommand{\NsfcHeader}[1]{%
    \stepcounter{NsfcHeaderCounter}%
    \setcounter{section}{0}%
    % 【关键】遇到全局大标题，把紧接在后面的所有子级标题顶部间距全部强行归零
    \global\setlength{\SecTopCurrent}{0pt}%
    \global\setlength{\SubSecTopCurrent}{0pt}%
    \par
    \vspace{\NsfcHeaderTopSkip}
    \noindent\hspace{\NsfcHeaderIndent}
    {\color{NsfcBlue}\sihao\mykai #1}
    \par
    \vspace{0.0em}
    \global\IsFirstSubHeadertrue
}

\newcommand{\NsfcSubHeader}[1]{%
    \par
    \ifIsFirstSubHeader
        \vspace{0.0em}
        \global\IsFirstSubHeaderfalse
    \else
        \vspace{0.2em}
    \fi
    \noindent\hspace{\NsfcSubHeaderIndent}
    {\sihao\color{NsfcBlue}\mykai\ziju{-0.00} #1}
    \vspace{0mm}
    % 【关键】遇到 NSFC 二级大标题，同样杀掉紧跟的正文标题间距
    \global\setlength{\SecTopCurrent}{0pt}%
    \global\setlength{\SubSecTopCurrent}{0pt}%
}

% --- 系统各级标题排版设置 (\section 等) ---
\ifcase\NumberingStyle
    % Case 0: (空)
\or % Case 1: 标准
    \renewcommand{\thesection}{\arabic{section}}
    \renewcommand{\thesubsection}{\arabic{section}.\arabic{subsection}}
    \renewcommand{\thesubsubsection}{\arabic{section}.\arabic{subsection}.\arabic{subsubsection}}
\or % Case 2: Word点号
    \renewcommand{\thesection}{\arabic{section}.}
    \renewcommand{\thesubsection}{\arabic{section}.\arabic{subsection}}
    \renewcommand{\thesubsubsection}{\arabic{section}.\arabic{subsection}.\arabic{subsubsection}}
\or % Case 3: 中文
    \renewcommand{\thesection}{（\chinese{section}）} 
    \renewcommand{\thesubsection}{\arabic{subsection}.}
    \renewcommand{\thesubsubsection}{（\arabic{subsubsection}）}
\fi

% --- 动态间距恢复机制 ---
% 逻辑：在排版标题主体之前（此时上方间距已经按照 0pt 结算完毕），将下一个同级标题的间距恢复为常态，并将子级标题的间距归零。
\titleformat{\section}{\zihao{4}\HeadingFont\bfseries}{\thesection}{0.5em}{%
    \global\setlength{\SecTopCurrent}{\SecTopBase}% 恢复下一个 \section 的间距
    \global\setlength{\SubSecTopCurrent}{0em}% 杀掉当前 \section 下第一个 \subsection 的间距
}
\titlespacing{\section}{0pt}{\SecTopCurrent}{1.0ex plus .2ex minus .2ex}

\titleformat{\subsection}{\zihao{4}\HeadingFont\bfseries}{\thesubsection}{0.5em}{%
    \global\setlength{\SubSecTopCurrent}{\SubSecTopBase}% 恢复下一个 \subsection 的间距
    \global\setlength{\SubSubSecTopCurrent}{0em}% 杀掉当前 \subsection 下第一个 \subsubsection 的间距
}
\titlespacing{\subsection}{0pt}{\SubSecTopCurrent}{0.5ex plus .1ex minus .2ex}

\titleformat{\subsubsection}{\zihao{-4}\HeadingFont\bfseries}{\thesubsubsection}{0.5em}{%
    \global\setlength{\SubSubSecTopCurrent}{\SubSubSecTopBase}% 恢复下一个 \subsubsection 的间距
}
\titlespacing{\subsubsection}{0pt}{\SubSubSecTopCurrent}{0.3ex plus .1ex minus .1ex}

% =========================================================
% 模块 5：核心文档状态注入 (合并全部 \AtBeginDocument)
% =========================================================

% 保留旧版废弃开关，防止外部主文档误调用导致直接报错
\newif\ifUseKaiBody
\UseKaiBodyfalse 

% 定义隐身模式开关
\newif\ifHideBody
\HideBodyfalse

% 统一管控正文字体、字号与隐身状态
\AtBeginDocument{
    \ifHideBody
        % 【隐身模式】: 字体变白，字号极小，行距压扁
        \color{white}
        \fontsize{1pt}{1pt}\selectfont
        \linespread{0.1}\selectfont
    \else
        % 【正常模式】: 强制正文小四号，并根据 style.tex 决定是否启用楷体
        \xiaosihao 
        \ifUseKaitiBody
            \ifdefined\mykai \mykai \else \kaishu \fi
        \else
            \ifdefined\mysong \mysong \else \songti \fi
        \fi
    \fi
}

% =========================================================
% 局部修复公式环境上下间距 (仅针对 equation 和 align 生效)
% =========================================================
\newcommand{\TightEquationSpacing}{%
    % 基础间距 0.5em，允许上下浮动 0.1em 到 0.2em 的弹性空间
    \setlength{\abovedisplayskip}{0.6em plus 0.2em minus 0.1em}%
    \setlength{\belowdisplayskip}{0.6em plus 0.2em minus 0.1em}%
    % 短行公式上方间距需进一步压缩，保持视觉紧凑
    \setlength{\abovedisplayshortskip}{0.2em plus 0.1em}%
    \setlength{\belowdisplayshortskip}{0.5em plus 0.2em minus 0.1em}%
}

\AtBeginEnvironment{equation}{\TightEquationSpacing}
\AtBeginEnvironment{align}{\TightEquationSpacing}
% 如果你还用了其他的 amsmath 环境，比如 gather, flalign 等，按同理追加即可

%--------------------------------------------------------
% 段间距
%--------------------------------------------------------
% 1. 声明长度变量
\newlength{\NsfcParskip}

% 2. 执行段间距设置 (利用 plus/minus 增加弹性)
\setlength{\NsfcParskip}{\UserParskip plus 0.2ex minus 0.1ex}
\setlength{\parskip}{\NsfcParskip}

% 3. 执行列表设置 (由于包含 if 逻辑，建议放在 AtBeginDocument 避错)
\AtBeginDocument{
    \ifnum\UserListMode=0
        \setlist{nosep} % 彻底无间距
    \else
        \setlist{
            partopsep=0pt,
            topsep=0.5\NsfcParskip,
            parsep=\NsfcParskip,
            itemsep=0.3\NsfcParskip
        }
    \fi
}

% 2. 独立设置页边距
\usepackage[a4paper, left=3cm,right=3cm,top=2.54cm,bottom=2.54cm]{geometry}

\begin{document}

\begin{center}
% 标题保持原样
{\sanhao \mykai \bfseries 报告正文（2026版）}
\end{center}
\vspace{-1mm} 

% 说明文字保持原样
{\sihao \mykai 参照以下提纲撰写，要求内容翔实、清晰，层次分明，标题突出。{\color{NsfcBlue} \bfseries 请勿删除或改动下述提纲标题及括号中的文字。}}
\vspace{-0.4em}
% =======================================================
% （一）立项依据
% =======================================================
\NsfcHeader{\textbf{（一）立项依据}}

\NsfcSubHeader{{\bfseries 科研仪器或核心部件研制的重要性和必要性} 论述研制的科研仪器或核心部件拟解决的重要科学和技术问题，同类科研仪器或核心部件的国内外研究现状和发展趋势。}



\section{一级标题示例}

标题的宏观风格由 \texttt{style.tex} 控制：
\begin{itemize}
  \item \textbf{换字体}：修改以下命令中的字体宏（可改为 \verb|\heiti| 或 \verb|\songti|）：
  \begin{codebox}
  \verb|\newcommand{\HeadingFont}{\kaishu}|
  \end{codebox}
  
  \item \textbf{换编号}：修改以下数字即可切换预设样式（1=极简理科；2=Word风；3=公文风）：
  \begin{codebox}
  \verb|\newcommand{\NumberingStyle}{2}|
  \end{codebox}
\end{itemize}

标题的细微排版由 \texttt{config.tex} 控制（尽量勿动）：搜索 \verb|\titleformat| 可更改更细致标题段落格式。

示例自定义加粗：\mybf{示例加粗},加粗自定义颜色，颜色可以在 \texttt{style.tex} 中调
\begin{itemize}
    \begin{codebox}
    \verb|\mybf{这里是加粗内容}|
    \end{codebox}
\end{itemize}

如果需要引用参考文献，引用格式如下：单个文献引用的样式~\citep{Ref1}，两个非连续文献引用的样式~\cite{Ref1, Ref3}，多个连续文献自动合并的样式~\cite{Ref1, Ref2, Ref3}。

\subsection{二级标题示例}
二级标题。

\subsubsection{三级标题示例}
三级标题。

\section{第二个一级标题}

\subsection{第二个二级标题}
开启新章节。


\begin{spacing}{1.1}
\vspace{0.5em}
{\wuhao \rmfamily % 添加 \rmfamily 强制切回 Times New Roman
\bibliographystyle{gbt7714-nsfc}
\renewcommand{\refname}{参考文献\vspace{0em}}
\bibliography{ref}
\vspace{11bp}
}
\end{spacing}



% =======================================================
% （二）研制方案和计划
% =======================================================
\NsfcHeader{\textbf{（二）研制方案和计划}}

\NsfcSubHeader{{\bfseries 1．研制内容与方案}科研仪器或核心部件的设计思想、总体结构等\linebreak 
\hspace*{0.1em}（请在此重点描述原理创新、技术创新或独到之处）；技术性能与主要\linebreak 
技术指标（含主要技术指标的提出依据和科学意义）；技术路线及设计\linebreak 
图；关键核心技术和解决方案；可购置或集成的部分以及配套部件解\linebreak 
决方案。}



这里写研制内容与方案。

\NsfcSubHeader{{\bfseries 2．拟解决的关键科学问题，拟取得的具有自主知识产权的关键技术}仪器研制过程中拟解决的关键科学或技术难题，拟取得的具有自主知识产权的关键技术；完成研制后，拟使用该仪器或核心部件解决的关键科学或技术问题。}

这里写拟解决的关键科学问题。

\NsfcSubHeader{{\bfseries 3.\hspace{0.0em}研制方案可行性分析}\hspace{-0.05em}（含研制风险与不确定性分析、应对措施{\bfseries）}。}
\vspace{-1.8em}

这里写可行性分析。

\NsfcSubHeader{{\bfseries 4．预期成果。}}

这里写预期成果。

\NsfcSubHeader{{\bfseries 5.\hspace{0.1em}验收指标的可考核性。}}

这里写验收指标。

\NsfcSubHeader{{\bfseries 6．年度研制计划}（含项目中期检查的阶段性目标）。}

这里写年度研制计划。


% =======================================================
% （三）研制基础及条件保障
% =======================================================
\NsfcHeader{{\bfseries （三）研制基础及条件保障}}

% 【修正】删除了原来这行下面的 \vskip 1mm
\NsfcSubHeader{{\bfseries 1．研制基础} 理论基础、技术基础；国内外可供利用的技术资源；人才队伍状况。}

这里写研制基础。

% 【修正】删除了原来这行下面的 \vskip 1mm
\NsfcSubHeader{{\bfseries 2．工作条件} 实验场地、实验环境、公用配套设施、研究人员的 \linebreak 时间保证等。}

这里写工作条件。

\NsfcSubHeader{{\bfseries 3.\hspace{0.1em}项目实施管理和保障措施}（申请人组织实施能力，项目运行管理模式，知识产权管理措施，项目档案与固定资产管理措施，依托单 \linebreak 位的配套措施与政策支持）。}

这里写管理和保障措施。

\NsfcSubHeader{{\bfseries 4.\hspace{0.1em}申请人和主要参与者仪器设备研制工作经验}[\hspace{-0.05em}其中，\hspace{-0.2em}如果申请人或主要参与者曾经或正在主持仪器设备研制或开发类项目，须在此列明项目的资助机构、项目类别、批准号、项目名称、获资助金额、起 \linebreak 止年月信息，并说明该项目的完成情况（已经结题的附结题报告摘要）以及该项目与本项目的关系][其中，如果申请人或主要参与者曾经主持部门推荐类国家重大科研仪器研制项目（含国家重大科研仪器设备研制专项项目）且已完成项目后评估的，须在此列明已完成后评估项\linebreak 目的后评估等级和后评估报告摘要]。}

这里写研制工作经验。


% =======================================================
% （四）其他需要说明的情况
% =======================================================
% 【修正】删除了原来这行上面的 \vskip 1mm
\NsfcHeader{{\bfseries（四）其他需要说明的情况}}

\NsfcSubHeader{{\bfseries 1.\hspace{0.3em}正在承担科研项目情况}（申请人和主要参与者正在承担的科研项目情况，要注明项目的资助机构、项目类别、批准号、项目名称、\linebreak 获资助金额、起止年月、与本项目的关系及负责的内容等；在“申请\linebreak 人和主要参与者仪器设备研制工作经验”中已说明的项目，无需重复\linebreak 录入）。}

无。

% 【修正】加上了 {\bfseries ...} 以保持风格统一
\NsfcSubHeader{2.\hspace{0.3em}申请人同年申请不同类型的国家自然科学基金项目情况（列明同年申请的其他项目的项目类型、项目名称信息，并说明与本项目之\linebreak 间的区别与联系；已收到自然科学基金委不予受理或不予资助决定的，无需列出）。}

无。

% 【修正】加上了 {\bfseries ...} 以保持风格统一
\NsfcSubHeader{3.\hspace{0.3em}具有高级专业技术职务（职称）的申请人或者主要参与者存在
\linebreak 同年申请或者参与申请国家自然科学基金项目的单位不一致的情况\linebreak \hspace*{0.0em}（列明所涉及人员的姓名，申请或参与申请的其他项目的项目类型、\linebreak 项目名称、单位名称、上述人员在该项目中是申请人还是参与者，并说明单位不一致原因）。}

无。

% 【修正】加上了 {\bfseries ...} 以保持风格统一
\NsfcSubHeader{4.\hspace{0.3em}具有高级专业技术职务（职称）的申请人或者主要参与者存在与正在承担的国家自然科学基金项目的单位不一致的情况（列明所涉及人员的姓名，正在承担项目的批准号、项目类型、项目名称、单位\linebreak 名称、起止年月，并说明单位不一致原因）。}

无。

% 【修正】加上了 {\bfseries ...} 以保持风格统一
\NsfcSubHeader{5.\hspace{0em}申请人和主要参与者同年以不同专业技术职务（职称）申请或参与申请科学基金项目的情况（应详细说明原因）。}

无。

% 【修正】将原本硬编码的第6点改为 \NsfcSubHeader，并加上了 {\bfseries ...}
\NsfcSubHeader{6.\hspace{0.3em}其他。}

无。

\end{document}